\section{1794? Adrien-Marie Legendre} % lang-pl
\label{sekcja_legendre}

Trzecią osobą, która zaatakuje problem prawdziwości piątego postulatu przez sprowadzenie do niedorzeczności jego alternatyw, będzie Adrien-Marie Legendre.
Rozpatrzy trzy możliwości: suma kątów trójkąta jest mniejsza, równa lub większa od $\pi$.
Tak jak Saccheri i Lambert założy milcząco, że proste są nieskończone i wyeliminuje pierwszą możliwość, ale nie poradzi sobie z trzecią.
Prosty styl, w jakim napisane będą dowody Legendre'a, oraz jego wysoka pozycja w świecie matematyki wzbudzą znaczne zainteresowanie problemem wśród innych osób.

EVES \cite[s. 286, 290]{eves1_1972}

\color{red}
Wyniki zasadniczo równoważne wynikom Saccheriego odkryje pół wieku później niezależnie od niego; jego \emph{Éléments de géométrie} (1794) doczekają się wielu wydań, przedruków i tłumaczeń, które ożywią zainteresowanie problemem równoległych \cite[s. 306]{hartshorne_2000}. % lang-pl
Jego \emph{Éléments de géométrie} (1794) przez około sto lat będą wiodącym podręcznikiem elementarnej geometrii: uprości i przestawi twierdzenia Euklidesa, a dwanaście wydań ukaże się do 1823 roku. % lang-pl
Pierwszy angielski przekład wykona w 1819 roku John Farrar z Harvardu, drugi w 1824 roku Thomas Carlyle (33 amerykańskie wydania), i tak podręcznik stanie się wzorem dla amerykańskich książek do geometrii \cite[s. 286]{eves1_1972}. % lang-pl
\index[persons]{Legendre, Adrien-Marie}%
\index[persons]{Farrar, John}%
\index[persons]{Carlyle, Thomas}%

Legendre zacznie od nowa: rozważy trzy hipotezy zależnie od tego, czy suma kątów trójkąta jest mniejsza od, równa albo większa od dwóch kątów prostych. % lang-pl
Przy cichym założeniu nieskończoności prostej wyeliminuje trzecią, ale mimo powtarzanych prób nie poradzi sobie z pierwszą; kolejne próby ukażą się w kolejnych wydaniach podręcznika, a ostatnią opublikuje w 1833 roku, w roku śmierci \cite[s. 286]{eves1_1972}. % lang-pl

\begin{theorem}
[pierwsze twierdzenie Legendre'a] % lang-pl
\index{twierdzenie!Legendre'a (o sumie kątów)} % lang-pl
	Suma kątów trójkąta nie jest większa od dwóch kątów prostych. % lang-pl
\end{theorem}

Dowód (przy założeniu nieskończoności prostej i aksjomatu Archimedesa): Eves \cite[s. 290, zad. 11]{eves1_1972}; Hartshorne \cite[s. 320--321]{hartshorne_2000} (tw. 35.2, twierdzenie Saccheriego--Legendre'a, w płaszczyźnie Hilberta z aksjomatem Archimedesa). % lang-pl
Hartshorne wybiera metodę Legendre'a, bo dowód Saccheriego korzysta z argumentu ciągłościowego; w obu przypadkach dowód istotnie używa aksjomatu Archimedesa dla kątów (lemat 35.1) \cite[s. 319]{hartshorne_2000}. % lang-pl
Greenberg \cite[s. 206, 209]{greenberg_2010} przypisuje ten wynik Saccheriemu i Legendre'owi (zwykle nazywa się go twierdzeniem Saccheriego--Legendre'a) i dodaje, że wystarcza słabszy od aksjomatu Archimedesa aksjomat Arystotelesa (zobacz sekcję \ref{sekcja_dobre_stare}); wtedy wyklucza się suma kątów większa od $\pi$. % lang-pl

\begin{theorem}
[drugie twierdzenie Legendre'a] % lang-pl
	Jeśli istnieje jeden trójkąt o sumie kątów równej dwóm kątom prostym, to suma kątów każdego trójkąta jest równa dwóm kątom prostym. % lang-pl
\end{theorem}

Dowód: Eves \cite[s. 290--291, zad. 13]{eves1_1972}; tam też wersja dla sumy mniejszej niż $\pi$. % lang-pl

Hartshorne \cite[s. 322--323]{hartshorne_2000} pokazuje, jak Legendre będzie ,,dowodził'' postulatu: założy, że przez punkt wewnątrz kąta zawsze przechodzi prosta przecinająca oba jego ramiona (zobacz aksjomat Legendre'a w sekcji \ref{sekcja_dobre_stare}). % lang-pl
Gdyby tak nie było, to dla dowolnie małego kąta istniałaby prosta leżąca w całości wewnątrz kąta (zad. 35.4), a to -- napisze Legendre -- „jest sprzeczne z naturą linii prostej'' (fr.~\emph{il répugne à la nature de la ligne droite}); wydrukuje to rozumowanie m.in. w swoim podręczniku z 1823 roku \cite[s. 322--323]{hartshorne_2000}. % lang-pl
Trzeciego twierdzenia Legendre nie udowodni, bo ono nie jest prawdziwe: nie da się wykluczyć sumy kątów mniejszej niż $\pi$. % lang-pl
W jednym z podejść zechce zbudować trójkąt zawierający dany trójkąt co najmniej dwa razy i w tym celu założy, że przez punkt wewnątrz kąta mniejszego niż $60^\circ$ zawsze da się poprowadzić prostą przecinającą oba ramiona -- czyli coś równoważnego piątemu postulatowi (Eves \cite[s. 290, zad. 12]{eves1_1972}). % lang-pl

Jego proste i przejrzyste dowody będą tak rozpowszechnione i on sam tak znany, że w całej Europie wzbudzi zainteresowanie problemem równoległych. % lang-pl
Rekord wytrwałości w próbach dowodu postulatu należy prawdopodobnie do niego \cite[s. 286]{eves1_1972}. % lang-pl
Nic dziwnego, że sprzeczności nie znajdzie: geometria z hipotezy kąta ostrego jest równie niesprzeczna jak euklidesowa \cite[s. 286]{eves1_1972}. % lang-pl
% Źródła: Hartshorne s. 304--311, 318 (§34, zad. 34.13--34.14); Greenberg 2010, s. 208.

\color{black}